\begin{circuitikz}[scale=0.9, transform shape]
  % Terminal de entrada
  \draw (0,0) node[ground] {}
    to[open, v^>=$v_i$] (0,2.5)
    to[short, o-*] (1.5,2.5) coordinate (in_node);

  % Resistencia RB = R1 // R2
  \draw (in_node) to[R, l=$R_B$, -] (1.5,0) node[ground] {};

  % Modelo híbrido del transistor
  \draw (in_node) to[short] (3.0,2.5)
    to[R, l=$h_{ie}$, i=$i_b$] (3.0,0) node[ground] {};

  % Fuente dependiente de corriente de colector
  \draw (5.5,0) node[ground] {}
    to[cisourceAM, l=$h_{fe} i_b$] (5.5,2.5)
    to[short, -] (7.0,2.5) coordinate (c_node);

  % Resistencia de colector RC
  \draw (7.0,2.5) to[R, l=$R_C$, -] (7.0,0) node[ground] {};

  % Resistencia de carga RL
  \draw (c_node) to[short] (8.5,2.5) coordinate (l_node)
    to[R, l=$R_L$, -] (8.5,0) node[ground] {};

   % 1. Dibujamos el terminal de salida
  \draw (l_node) to[short, -o] (9.5,2.5) node[above] {$v_L$} coordinate (out_node);

  % 2. Colocamos la masa en el punto inferior y medimos la tensión en abierto hacia arriba
  \draw (9.5,0) node[ground] {} to[open, v_>=$v_L$] (out_node);
\end{circuitikz}
